\section{Illumination, speech, and the order of spirits}

\subsection{Question 106: Sharing the light of truth}\label{q106}
An angel's knowledge is personal without being a private possession held against the good of others. Aquinas begins his account of angelic society with illumination: one intellect helps another to understand. The higher angel strengthens the lower angel's intellectual activity and unfolds a truth in a manner proportioned to the recipient. A human teacher provides Aquinas's analogy. Understanding a subject as a whole, the teacher divides its explanation into steps that a learner can grasp. The angels do not need our laborious reasoning, but the proportion between teaching and capacity remains. This communication neither creates the recipient's intellect nor bestows grace or the light of glory. Those gifts come from God.\footnote{\ST{I q.~106 a.~1}, corpus and ad~2--3.}

Nor does an angel become the medium through which another sees God's essence. Every blessed angel sees God immediately. What one communicates to another concerns the divine works and their reasons: seeing God does not mean comprehending everything God knows. Thus illumination belongs within the communion of those already blessed, rather than filling a gap between a blessed creature and God.\footnote{\ST{I q.~106 a.~1 ad~1}.}

Article~2 draws a boundary around the will. An angel can disclose something lovable and so persuade another to love it. He cannot produce the other's act of willing from within. The power of willing belongs to a creature whose author is God; only God reaches it as its inward cause. Consequently, Dionysian purification does not mean that one blessed angel absolves another of sin. Here purification removes a lack of understanding; perfection brings understanding to its proper completion.\footnote{\ST{I q.~106 a.~2}, corpus and ad~1.}

In article~3, illumination proceeds downward through the heavenly order. Aquinas does not transfer this rule without qualification to ecclesiastical office: a person lower in office can possess greater holiness or knowledge and teach a superior. His account of heaven rests on the correspondence between spiritual excellence and hierarchical place which he does not suppose invariably present on earth.\footnote{\ST{I q.~106 a.~3}, especially ad~1.}

Gregory describes the heavenly city as a communion in which each order's special excellence becomes the good of all through charity. All burn with love, all drink from the source of knowledge, and all receive God's presence, though their names mark differing fullness of gift.\footnote{Gregory, \emph{Homilies on the Gospels} II.34.14.}

Article~4 joins generosity to enduring difference. Higher angels communicate the gifts they receive as fully as recipients can receive them; the difference lies in the manner of possessing the truth, not a refusal to share. A teacher and pupil can know the same conclusion with unequal depth. New divine disclosures concerning providence and salvation also give fresh matter for communication before the final judgment.\footnote{\ST{I q.~106 a.~4}, corpus and ad~1--3.}

\angeldossiersettled{I q.~106, aa.~1--4}{Thomist position: intellectual illumination and its communication.}{Dionysius, \emph{Celestial Hierarchy} VII--VIII, XV; Gregory, \emph{Homilies on the Gospels} II.34.14; the article also invokes Augustine and Peter Lombard.}

\subsection{Question 107: Speaking without sound}\label{q107}
Angelic speech is the voluntary communication of an intellectual conception. The absence of vocal organs removes the need for a sound; it does not remove the distinction between having a thought and addressing it to someone. Aquinas distinguishes knowledge possessed habitually, knowledge actually considered, and a conception intentionally directed to another. Speech occurs at the third stage. Created minds are therefore not automatically transparent to one another. God alone knows the inward thought without needing the creature to disclose it.\footnote{\ST{I q.~107 a.~1}, corpus and ad~1--2.}

This permits a lower angel to speak to a higher one. The statement that a speaker wishes to do something makes the speaker's intention known; it need not confer the higher grasp of divine truth which Aquinas calls illumination. Every illumination is speech in this extended sense, but not every speech is illumination. The hierarchy of teaching does not reduce the society of angels to a one-way stream of commands.\footnote{\ST{I q.~107 a.~2}.}

Angels also speak to God. They do not supply information to the omniscient Creator. They turn toward him in admiration and in consultation about their ministry. Aquinas thus distinguishes communication intended to instruct the hearer from communication intended to receive instruction. Praise is continual; consultation about a newly assigned work answers a particular occasion.\footnote{\ST{I q.~107 a.~3}, corpus and ad~2.}

Distance does not weaken angelic speech, because it is an intellectual act rather than a sound carried through a medium. This does not make the speaker omnipresent. The operation belongs to the speaker where he is; its intelligibility is not measured by the distance to the hearer. Nor must every angel hear every communication. Its recipients are determined by the speaker's will, not by the reach of a voice. Articles~4 and~5 therefore deny both a physical transmission barrier and an automatic universal audience.\footnote{\ST{I q.~107 aa.~4--5}.}

\angeldossiersettled{I q.~107, aa.~1--5}{Thomist position: voluntary intellectual communication.}{1~Corinthians 13:1; Zacharias 1:12; Luke 16:24; Dionysius, \emph{Celestial Hierarchy} VII.3. Aquinas also invokes Gregory's \emph{Moralia}.}

\subsection{Question 108: One communion, several orders}\label{q108}
The word \emph{hierarchy} names a sacred ordering toward God. Article~1 can therefore speak of one society under one divine ruler, while distinguishing several ways of receiving and communicating his gifts. Humans understand through sensible signs; angels receive intellectual illumination without that dependence. Within angelic knowledge, Aquinas distinguishes attention to the divine reasons of things in God, to their universal created causes, and to their application in particular effects. These correspond to the three hierarchies. The distinction does not place the Trinity at the top of a ladder of unequal gods: Aquinas expressly rejects a hierarchy within the divine Persons.\footnote{\ST{I q.~108 a.~1}.}

Article~2 distinguishes orders within each hierarchy through their offices. The threefold arrangement marks beginning, middle, and completion within an ordered whole. Article~3 prevents this classification from becoming an exhaustive celestial census. Many angels belong to one order because our knowledge grasps their offices generally. Each angel has a more particular place and office than the nine names tell us. Membership in a common order does not imply exact equality of natural or gracious gifts.\footnote{\ST{I q.~108 aa.~2--3}.}

Nature and grace both matter in article~4. A natural order can be considered with respect to the knowledge and love of God accessible to created powers. The order consummated in the vision of God depends on grace. Aquinas holds that angels receive gratuitous gifts proportioned to their natural capacities; he expressly refuses to make that proportionality a universal rule for human beings. Natural brilliance is therefore no scale for measuring a person's sanctity.\footnote{\ST{I q.~108 a.~4}; compare I q.~62 a.~6.}

The names of the orders, treated in article~5, signify characteristic excellence rather than exclusive possession. All holy angels love and know God; the Seraphim are named from an eminent ardor of love, the Cherubim from fullness of knowledge. The higher possess the perfections of the lower more eminently, while the lower participate in higher perfections according to capacity. ``Angels'' can consequently name all the heavenly messengers or the lowest order in particular. Likewise, ``virtue'' can signify power generally or the specific order called Virtues.\footnote{\ST{I q.~108 a.~5}, especially ad~1 and ad~5--6.}

Article~6 compares two received sequences. In the customary Latin terminology, Dionysius's descending arrangement is Seraphim, Cherubim, Thrones; Dominations, Virtues, Powers; Principalities, Archangels, Angels. Gregory's descending list is Seraphim, Cherubim, Thrones, Dominations, Principalities, Powers, Virtues, Archangels, Angels. Thus Principalities and Virtues exchange positions, while Powers retain their place between them. Parker's English ``Powers'' for the Dionysian \emph{dynameis} corresponds here to Latin \emph{Virtutes}; his ``Authorities'' corresponds to \emph{Potestates}. An English word alone cannot settle the comparison.\footnote{Dionysius, \emph{Celestial Hierarchy} VI--IX, trans. John Parker, \emph{Works}, vol.~II (1899), pp.~23--38; Gregory, \emph{Homilies on the Gospels} II.34.7--10, Latin electronic witness; \ST{I q.~108 a.~6}.}

Aquinas considers both arrangements intelligible. Dionysius emphasizes the reception and transmission of divine illumination; Gregory's explanations emphasize ministries: Principalities preside over good spirits, Powers restrain hostile spirits, and Virtues minister in wonders. Aquinas argues that the functions assigned under different names can be close in substance. His reconciliation does not make their written sequences identical. Nor does a list establish that every member of an order is engaged only in the one activity from which its name comes.\footnote{\ST{I q.~108 aa.~5--6}, especially a.~6 ad~4; Gregory, \emph{Homilies} II.34.10,14.}

In article~7 the differences of nature and glory remain after judgment, while ministries directed to bringing pilgrims to their goal cease in that form. Communion and illumination are not abolished when salvation is complete. Aquinas compares the distinction between an army's work in battle and its order in triumph. Article~8 then places redeemed humanity within the one blessed society. Grace can raise human beings to a glory equal to that of angels without changing human nature into angelic nature. The claim concerns participation in blessedness, not a transformation of a departed person into a different species of creature. Aquinas also rejects the restriction of this shared society to virgins or the specially perfect, as though other saved persons formed a separate society.\footnote{\ST{I q.~108 aa.~7--8}.}

\angeldossier{I q.~108, aa.~1--8}{Thomist position: hierarchy and participation. Disputed: the arrangement of Principalities and Virtues.}{Dionysius, \emph{Celestial Hierarchy} VI--IX; Gregory, \emph{Homilies} II.34.7--14; Ephesians 1:20--21; Colossians 1:16; Isaias 6:2--3; Matthew 22:30.}{Gregory places Principalities where Dionysius places Virtues, and Virtues where Dionysius places Principalities. Aquinas's functional reconciliation accompanies, rather than erases, this difference.}

\subsection{Question 109: Order surviving rebellion}\label{q109}
The demons retain the natural powers of spiritual creatures. Article~1 therefore distinguishes natural order, the grace in which the angels were created according to Aquinas, and consummated glory. Fallen angels retain the first, lost the second, and never possessed the third. Their continuing powers are goods of a created nature; wickedness is their misuse. Names whose meanings imply perfected charity or divine indwelling cannot simply be assigned to demons because of a supposed former rank.\footnote{\ST{I q.~109 a.~1}, especially ad~3.}

Article~2 finds precedence among demons in their remaining inequalities of nature. Their cooperation does not imply charity or friendship. It can arise from a shared purpose of harming human beings. Leadership in such activity increases misery rather than conferring a genuine blessed excellence. Article~3 likewise distinguishes communication from illumination: demons can make their intentions known, but do not illuminate one another in the proper sense of leading understanding toward God.\footnote{\ST{I q.~109 aa.~2--3}.}

Good angels exercise authority over bad angels because their adherence to divine justice is stronger than a demon's natural superiority. A naturally lower blessed angel can therefore restrain a naturally higher fallen spirit. Aquinas invokes the order Augustine describes, in which the sinful rational spirit is governed by the just spirit, and the just spirit by God. This government permits some harms under providence; it does not mean that the good angels endorse demonic malice or that evil has escaped their Lord's jurisdiction. Information imparted to demons in executing a divine judgment is received perversely by them, even when communicated for a righteous purpose.\footnote{\ST{I q.~109 a.~4}, corpus and ad~1--3; Augustine, \emph{De Trinitate} III.4.9; Gregory, \emph{Homilies} II.34.10.}

\angeldossiersettled{I q.~109, aa.~1--4}{Thomist position: the ordering of fallen spirits. Their subjection to providence belongs to Church teaching; see the Catechism, no.~395.}{Ephesians 6:12; Ecclesiasticus 34:4; Augustine, \emph{De Trinitate} III.4.9; Gregory, \emph{Homilies} II.34.10.}
